\documentclass[11pt]{article}

\usepackage[letterpaper,margin=1in]{geometry}
\usepackage{amsmath,amssymb,amsfonts}
\usepackage{booktabs}
\usepackage{graphicx}
\usepackage{hyperref}
\usepackage{xcolor}
\usepackage{enumitem}

\hypersetup{colorlinks=true,linkcolor=blue,urlcolor=blue}
\setlength{\parindent}{0pt}
\setlength{\parskip}{0.6em}

\newcommand{\studentname}{YOUR NAME}
\newcommand{\andrewid}{YOUR ANDREW ID}
\newcommand{\answerlabel}[2]{%
  \clearpage
  \section*{#1: #2}
  \addcontentsline{toc}{section}{#1: #2}
}
\newenvironment{response}{\color{black}}{}

\begin{document}

\begin{center}
  {\Large 16-720 Computer Vision, Fall 2026}\\[0.4em]
  {\LARGE Programming Assignment 3: Geometry in Motion}\\[1em]
  \textbf{Name:} \studentname \hfill \textbf{Andrew ID:} \andrewid
\end{center}

\tableofcontents

\answerlabel{A1}{Derive the DLT constraints}
\begin{response}
% Answer A1.1--A1.5 here. Keep every subpart visibly labeled.
\textbf{A1.1}\par
\textbf{A1.2}\par
\textbf{A1.3}\par
\textbf{A1.4}\par
\textbf{A1.5}\par
\end{response}

\answerlabel{A2}{Normalize a small correspondence set}
\begin{response}
\textbf{A2.1}\par
\textbf{A2.2}\par
\textbf{A2.3}\par
\textbf{A2.4}\par
\textbf{A2.5}\par
\end{response}

\answerlabel{A3}{RANSAC sampling without replacement}
\begin{response}
\textbf{A3.1}\par
\textbf{A3.2}\par
\textbf{A3.3}\par
\textbf{A3.4}\par
\end{response}

\answerlabel{A4}{Local motion with additive brightness change}
\begin{response}
\textbf{A4.1}\par
\textbf{A4.2}\par
\textbf{A4.3}\par
\textbf{A4.4}\par
\textbf{A4.5}\par
\end{response}

\answerlabel{B1}{Read the geometry pipeline}
\begin{response}
\textbf{B1.1}\par
\textbf{B1.2}\par
\textbf{B1.3}\par
\textbf{B1.4}\par
\textbf{B1.5}\par
\end{response}

\answerlabel{B2}{Normalized DLT evidence}
\begin{response}
Summarize implementation decisions, tests, and any numerical failure evidence. Reference notebook figures rather than copying code.
\end{response}

\answerlabel{B3}{RANSAC evaluation}
\begin{response}
Include the required outlier-fraction table, corner error, precision/recall, runtime, and interpretation.
\end{response}

\answerlabel{B4}{Coarse-to-fine Lucas--Kanade translation}
\begin{response}
Include pyramid/sign checks, level traces, plots, and the initialization-versus-brightness analysis.
\end{response}

\answerlabel{C1}{Planar target and application}
\begin{response}
Document the capture and include a raw still, model rationale, and predicted violation.
\end{response}

\answerlabel{C2}{System design and substantive improvement}
\begin{response}
Describe the baseline state machine, validity gates, improvement, and diagnostic outputs.
\end{response}

\answerlabel{C3}{AI-comprehension reflection}
\begin{response}
Write 200--300 words and answer all four prompts. If no generative AI was used, follow the alternative in the specification.
\end{response}

\answerlabel{D1}{Controlled robustness study}
\begin{response}
Include the capture table, prior predictions, manual-corner metrics, method comparisons, representative views, and explanation.
\end{response}

\answerlabel{D2}{When a good-looking warp is wrong}
\begin{response}
Include success/failure evidence and answer all three interpretation questions.
\end{response}

\answerlabel{D3}{Hero Result}
\begin{response}
Embed the static preview, identify the submitted artifact, and state the assumption and non-claim shown in it.
\end{response}

\answerlabel{E}{Optional extension}
\begin{response}
Delete this page if no extension is submitted. State whether you opt into the public showcase.
\end{response}

\answerlabel{Acknowledgments}{Collaboration, resources, data, weights, and AI use}
\begin{response}
List every collaborator and outside resource. Describe any generative-AI use in permitted sections. State ``none'' explicitly where appropriate.
\end{response}

\end{document}
